\documentclass[11pt]{article}
\usepackage[margin=1in]{geometry}
\usepackage[T1]{fontenc}
\usepackage{lmodern,microtype,hyperref,enumitem,amsmath,amssymb,xcolor,booktabs}
\usepackage[numbers,sort&compress]{natbib}
\usepackage{url}
\hypersetup{colorlinks=true,urlcolor=blue,citecolor=blue,linkcolor=blue}

\title{Q12 --- Deeper projects\\\large Why do models give confident answers even when wrong?}
\author{Deep Learning Systems 2026 --- AI Internals Module}
\date{}

\begin{document}
\maketitle

\section*{Framing the question}

Confident wrong answers are a symptom, not a single mechanism. They arise from
the interaction of several internal properties: modern neural networks are
\emph{poorly calibrated}, with confidence systematically exceeding accuracy
\citep{guo2017calibration}; maximum-likelihood training pushes probability mass
toward sharp distributions regardless of epistemic uncertainty; and
reinforcement learning from human feedback can reward assertive, fluent
responses over hedged ones, amplifying overconfidence and hallucination
\citep{ouyang2022training,ji2023survey}. A further source is that language
models can be ``calibrated'' at the token level yet still produce factually
wrong claims stated fluently, because fluency and factuality are governed by
different internal pathways \citep{kadavath2022know}. The projects below ask you
to measure, localize, and partially mitigate these phenomena, treating
confidence as an internal computation that can be probed rather than an opaque
model trait.

\section*{Project 1 --- Calibration cartography across tasks and scale}

\textbf{Goal.} Build an empirical map of how calibration error relates to model
scale, task type, and decoding temperature, and determine whether
overconfidence is uniform or concentrated in specific regimes.

\textbf{Approach.} Select at least three model sizes from an open family.
Evaluate on a mix of multiple-choice benchmarks (where token probabilities over
answer letters give a clean confidence signal) and short free-form QA (where you
elicit a verbalized confidence or use sequence likelihood). Compute expected
calibration error (ECE), Brier score, and reliability diagrams following
\citet{guo2017calibration}, and separately measure selective-prediction quality
via risk--coverage curves. Contrast pretrained checkpoints against
instruction-tuned ones to test whether alignment tuning degrades calibration
\citep{kadavath2022know,ouyang2022training}.

\textbf{Deliverables.}
\begin{itemize}[leftmargin=1.5em]
  \item Reliability diagrams and ECE/Brier tables broken down by model size,
  task type, and temperature.
  \item A risk--coverage analysis with at least two selection scores.
  \item A 4--6 page report stating where overconfidence is worst and whether it
  is driven by scale, tuning, or decoding, with reproducible code.
\end{itemize}

\section*{Project 2 --- Does the model ``know'' it is wrong internally?}

\textbf{Goal.} Test whether a model's hidden states encode a usable signal of
its own correctness that is stronger than its stated confidence.

\textbf{Approach.} Collect model answers labeled correct/incorrect on a
factual-QA set. Train lightweight linear probes on frozen hidden
representations at several layers to predict correctness, following the
self-knowledge paradigm of \citet{kadavath2022know} and probing methodology from
\citet{alain2017understanding}. Compare probe AUROC against the model's own
token-level or verbalized confidence. Investigate layer-wise emergence: does an
internal correctness signal appear before the final layer, and is it discarded
by the output head? Optionally relate findings to representation-level factuality
signals in the hallucination literature \citep{ji2023survey}.

\textbf{Deliverables.}
\begin{itemize}[leftmargin=1.5em]
  \item Layer-wise probe AUROC curves vs.\ baseline confidence.
  \item An analysis of whether the ``knows-it-is-wrong'' signal is linearly
  decodable and where it peaks.
  \item A short report discussing implications for why the output layer can be
  confident when intermediate states are not.
\end{itemize}

\section*{Project 3 --- Mitigation: temperature, ensembling, and verbalized uncertainty}

\textbf{Goal.} Compare post-hoc and prompting-based mitigations for
overconfidence and quantify their accuracy/calibration trade-offs.

\textbf{Approach.} Implement three mitigations on a shared evaluation suite:
(a) temperature scaling fit on a held-out set \citep{guo2017calibration};
(b) deep-ensemble or sampling-based uncertainty, aggregating multiple stochastic
decodes and using self-consistency as a confidence proxy
\citep{lakshminarayanan2017simple,wang2023selfconsistency}; and (c) explicit
prompting for calibrated verbalized confidence. Measure ECE, selective accuracy,
and the cost in latency/compute for each.

\textbf{Deliverables.}
\begin{itemize}[leftmargin=1.5em]
  \item A unified benchmark harness applying all three mitigations.
  \item Before/after reliability diagrams and a cost--benefit table.
  \item A recommendation note on which mitigation suits latency-constrained vs.\
  accuracy-critical deployments.
\end{itemize}

\section*{Project 4 --- Training dynamics of overconfidence}

\textbf{Goal.} Trace how confidence miscalibration emerges during training or
fine-tuning, connecting it to loss-surface and memorization phenomena.

\textbf{Approach.} Fine-tune a small model while checkpointing frequently. At
each checkpoint, measure accuracy and ECE, and track whether the model becomes
confident on examples it later gets wrong. Relate observations to the tendency of
networks to fit noisy or memorized data with high confidence
\citep{zhang2021understanding} and to over-sharpening under cross-entropy. If
using preference fine-tuning, test whether reward optimization increases
assertiveness independent of correctness \citep{ouyang2022training}.

\textbf{Deliverables.}
\begin{itemize}[leftmargin=1.5em]
  \item Trajectory plots of accuracy and ECE across training steps.
  \item An analysis isolating when confidence decouples from correctness.
  \item A report linking the empirical dynamics to at least one theoretical
  account of overconfidence.
\end{itemize}

\section*{Submission notes}

Each project must ship with runnable code, fixed random seeds, a data manifest,
and a self-contained report. Grading weights reproducibility and the clarity of
the causal story linking an \emph{internal} mechanism to the observed symptom of
confident errors, not benchmark leaderboard position.

\begin{thebibliography}{9}

\bibitem[Alain and Bengio(2017)]{alain2017understanding}
G.~Alain and Y.~Bengio.
\newblock Understanding intermediate layers using linear classifier probes.
\newblock \emph{International Conference on Learning Representations (Workshop)},
2017.

\bibitem[Guo et~al.(2017)]{guo2017calibration}
C.~Guo, G.~Pleiss, Y.~Sun, and K.~Q. Weinberger.
\newblock On calibration of modern neural networks.
\newblock In \emph{International Conference on Machine Learning (ICML)}, 2017.

\bibitem[Ji et~al.(2023)]{ji2023survey}
Z.~Ji, N.~Lee, R.~Frieske, T.~Yu, D.~Su, Y.~Xu, E.~Ishii, Y.~J. Bang,
  A.~Madotto, and P.~Fung.
\newblock Survey of hallucination in natural language generation.
\newblock \emph{ACM Computing Surveys}, 55(12):1--38, 2023.

\bibitem[Kadavath et~al.(2022)]{kadavath2022know}
S.~Kadavath, T.~Conerly, A.~Askell, T.~Henighan, D.~Drain, E.~Perez,
  N.~Schiefer, Z.~Hatfield-Dodds, N.~DasSarma, et~al.
\newblock Language models (mostly) know what they know.
\newblock \emph{arXiv preprint arXiv:2207.05221}, 2022.

\bibitem[Lakshminarayanan et~al.(2017)]{lakshminarayanan2017simple}
B.~Lakshminarayanan, A.~Pritzel, and C.~Blundell.
\newblock Simple and scalable predictive uncertainty estimation using deep
  ensembles.
\newblock In \emph{Advances in Neural Information Processing Systems
  (NeurIPS)}, 2017.

\bibitem[Ouyang et~al.(2022)]{ouyang2022training}
L.~Ouyang, J.~Wu, X.~Jiang, D.~Almeida, C.~L. Wainwright, P.~Mishkin, C.~Zhang,
  S.~Agarwal, K.~Slama, A.~Ray, et~al.
\newblock Training language models to follow instructions with human feedback.
\newblock In \emph{Advances in Neural Information Processing Systems
  (NeurIPS)}, 2022.

\bibitem[Wang et~al.(2023)]{wang2023selfconsistency}
X.~Wang, J.~Wei, D.~Schuurmans, Q.~Le, E.~Chi, S.~Narang, A.~Chowdhery, and
  D.~Zhou.
\newblock Self-consistency improves chain of thought reasoning in language
  models.
\newblock In \emph{International Conference on Learning Representations (ICLR)},
2023.

\bibitem[Zhang et~al.(2021)]{zhang2021understanding}
C.~Zhang, S.~Bengio, M.~Hardt, B.~Recht, and O.~Vinyals.
\newblock Understanding deep learning (still) requires rethinking
  generalization.
\newblock \emph{Communications of the ACM}, 64(3):107--115, 2021.

\end{thebibliography}

\end{document}